\documentclass{article}

% ---- Packages ----
\usepackage{array}        % for p{width} columns
\usepackage{geometry}     % adjust margins
\usepackage{caption}      % better control of captions
\usepackage{booktabs}     % optional: nicer rules
\usepackage{longtable}    % optional: multipage tables
\usepackage{graphicx}     % in case you want to add diagrams

\geometry{margin=1in}

\begin{document}
\begin{table}[htbp]
\centering
\small % or \scriptsize for even smaller
\setlength{\tabcolsep}{4pt} % reduce horizontal padding
\renewcommand{\arraystretch}{1.1} % reduce vertical spacing
\caption{Comparison of Lesion Analysis and Mapping Tools}
\begin{tabular}{|p{2.8cm}|p{2.8cm}|p{2.8cm}|p{2.8cm}|p{2.8cm}|}
\hline
\textbf{Tool} & \textbf{Input} & \textbf{Output} & \textbf{Pros} & \textbf{Cons} \\
\hline
VLSM & Lesion masks + behavioral scores & Lesion–symptom localization & Well-established; anatomy–behavior link & Needs large number of scans; sensitive to registration \\
\hline
LQT (MATLAB) & MRI lesion masks; atlas & Estimates region-level grey matter damage and white matter disconnection & Publicly available & Atlas accuracy  \\
\hline
SVR-LSM & Lesion masks + behavior & ML model & Flexible; handles high-dimensional data & Needs behavioral data \\
\hline
DSM & Lesion masks + tract atlas & White matter disconnection maps & Highlights network disruptions & Requires diffusion atlas  \\
\hline
CLSM & Lesion masks + connectome & Network lesion mapping & Network-level insight; connectomics & Needs high-quality connectome \\
\hline
\end{tabular}
\end{table}

\end{document}
